
% ----------------------------------------------------------------------
%           Packages 
% ----------------------------------------------------------------------
\usepackage{lipsum}%dummy tekst
%\usepackage{todonotes}%todo's : https://tug.ctan.org/macros/latex/contrib/todonotes/todonotes.pdf

% -- layout hoofdstukken: Sonny, Lenny, Glenn, Conny, Rejne, Bjarne, Bjornstrup
%\usepackage[Glenn]{fncychap}
% -- andere manier voor hoofdstukken, maar misschien niet met hyperref : https://markov.htwsaar.de/tex-archive/macros/latex/contrib/titlesec/titlesec.pdf
\usepackage{xpatch}
\usepackage{titlesec}
\makeatletter

% \addcounter{chapter}

\makeatletter
\@ifundefined{c@chapter}
  {\newcounter{chapter}}
  {}
\makeatother

\titleformat{\chapter}[frame]
  {\normalfont}{\filright\enspace \@chapapp~\thechapter\enspace}
  {8pt}{\LARGE\bfseries\filcenter}
\titlespacing*{\chapter}
  {0pt}{0pt}{20pt}
%\@addtoreset{chapter}{part}. %-- als je de hoofdstuk nummering altijd opnieuw wilt hebben
% ---- met de 2 lijnen hieronder kan je onder een part title toch wat schrijven.
\xpatchcmd{\@endpart}{\vfil\newpage}{}{}{}%\vspace{10mm}
\xpatchcmd{\@endpart}{\newpage}{}{}{}

\makeatother


%-- QR codes
\usepackage{hyperref}
\usepackage{qrcode}
% -- language: dutch --
\usepackage{csquotes}
\usepackage[dutch]{babel}
% -- notes in margin
%Why use both marginnotes and sidenotes? Quite simply, marginnotes overlap
%each other if they are too close whereas sidenotes both numbers and dynamically aligns
%all side notes, figures, and tables. However sidenotes cannot be used in equations,
%multicols, and with the tcolorbox4 environment. As the majority of the special
%environments from amsthm are modified to use tcolorbox, marginnotes becomes an
%essential part of NotesTeX.
\usepackage{marginnote}
%\usepackage{sidenotes}
\usepackage{nuc}	%Nuclides  https://mirror.ibcp.fr/pub/CTAN/macros/latex/contrib/nuc/nuc.pdf 
\usepackage{ifthen}
% nodig voor sitenotes
\usepackage{xparse}
\usepackage{l3keys2e}
%\usepackage{changepage}
% -- nodig voor env en andere zaken
\usepackage{fancyhdr,geometry}
% \usepackage[usenames,dvipsnames]{xcolor}
 %\usepackage[many]{tcolorbox}
%\usepackage{varwidth}
\tcbuselibrary{raster}
\usepackage[T1]{fontenc}                            % Font Styling
\usepackage{lmodern,mathrsfs}
% \usepackage[shortlabels]{enumitem}                  % Enumitem Options
\usepackage{mathtools,amsthm,amssymb,amsfonts,bm}   % Math Presets
\usepackage{thmtools,amsmath}
\usepackage{array,tabularx,booktabs}                % Table Presets
\usepackage{graphicx,wrapfig,float,caption}         % Figure Presets
\usepackage[export]{adjustbox}
\usepackage{setspace,multicol}                      % Text Presets
\usepackage{tikz,physics}                           % Physics Presets
\usepackage{tkz-euclide}
\usetikzlibrary{arrows.meta, positioning}
\usepackage{circuitikz}
\usepackage{xcolor}						%text kleur
\usepackage{colortbl}
\usepackage{tabularx}
\colorlet{tableheadcolor}{gray!25} %eerste rij tabel grijs
\usepackage{caption}
\usepackage{subcaption}
\usepackage{cancel} %cancel streepjes
\usepackage{setspace} %voor interlinie: \onehalfspacing en \singlespacing
% -- file packages
\usepackage{pdfpages} %om pdf's toe te voegen
%\usepackage{filecontents}
%\usepackage{subfiles} % Best loaded last in the preamble
\usepackage{siunitx} % SI units en aanpassingen
\usepackage[locale=FR]{siunitx}
\AtBeginDocument{\RenewCommandCopy\qty\SI}
\usepackage{physics}
	\sisetup{exponent-product = \cdot, output-product = \cdot}
	\sisetup{per-mode = fraction}
	\sisetup{group-separator = {\,}} %kleine spatie tussen 1000 tallen
	\DeclareSIUnit\seconde{s}
	\DeclareSIUnit\jaar{j}
	\DeclareSIUnit\dag{d}
	\DeclareSIUnit\graden{^\circ C}
	\DeclareSIUnit\meter{m}
	\DeclareSIUnit\newton{N}
	\DeclareSIUnit\radiaal{rad}	
	\sisetup{inter-unit-product = \ensuremath{{}\cdot{}}} %nog nodig?
% % Margins
% %\oddsidemargin = 0pt
% %\marginparwidth = 85pt
% % dit werkt------------------------------
% \geometry{paperheight=11in,paperwidth=8.5in,
% marginparsep=.019\paperwidth,marginparwidth=.2\paperwidth,
% inner=.12\paperwidth,voffset=0in,headheight=.02\paperheight,hoffset=0.2in,
% headsep=.02\paperheight,footskip=20pt,
% textheight=.795\paperheight,textwidth=.62\paperwidth}
% %--------------------------------------------
\setcounter{tocdepth}{1} %om subsections niet in inhoudstafel te krijgen


% tekening ladingen
% tikz
\tikzstyle{charge}=[thin,top color=red!50,bottom color=red!70,shading angle=20]
\tikzstyle{charge+}=[thin,top color=red!50,bottom color=red!90!black,shading angle=20]
\tikzstyle{charge-}=[thin,top color=blue!50,bottom color=blue!80,shading angle=20]
\tikzstyle{force}=[->,very thick,orange!80!black]
\tikzstyle{vector}=[->,very thick,green!45!black]
\def\a{1.4}
\def\R{0.15}
\def\F{0.8}


% % --------
% % nummering aanpassen
% \usepackage{chngcntr}

% % 1. Reset chapters na een part
% \counterwithin{chapter}{part}

% % 2. Maak de opmaak van het chapter als Roman.Arabic (bijv. II.1)
% \renewcommand{\thechapter}{\Roman{part}.\arabic{chapter}}

% % 3. Zorg dat secties resetten na een chapter en toon als II.1.1
% \counterwithin{section}{chapter}
% \renewcommand{\thesection}{\thechapter.\arabic{section}}
% %-----------------------


\newlength{\spacelength}
\settowidth{\spacelength}{\normalfont\ }
\declaretheoremstyle[
    headfont={\bfseries\sffamily\footnotesize},
    notebraces={}{}, % <--- Verwijdert de openings- en sluithaakjes
	notefont={\bfseries},
    bodyfont={\normalfont},
    headpunct={\relax},%\newline,
    headformat={%
        \makebox[0pt][r]{\NAME\ \NUMBER\hspace{\marginparsep}\hspace{1cm}}\hskip-\spacelength{\normalsize\NOTE}},
]{definition}


% Logo command
\ifdefined\isXourse
\logo{xmPictures/buisjes.jpg}
\else

% \logo is only defined in xourses (but that prevents \logo to be used in global.sty etc ....)

% \renewcommand*{\logo}[1]{%
%     \IfFileExists{#1}%
%       {\gdef\xourse@logo{#1}}%
%       {\ClassError{xmPreamble}{logo file does not exist}
%         {To use logo, make sure that the referenced image file exists}}%
% }

 \fi


 % Mylink: doet nog niet wat ik dacht dat het doet!!!
 \newcommand{\mylink}[2][]{
	\ifdefined\HCode %Href werkt niet?
		\nl\url{#2}: \text{#1} 
	\else

		\qrcode[hyperlink]{#2}#1
		%figure werkt niet goed.
		% \begin{figure}[h!]
		% 	\qrcode[hyperlink]{#2}
		% 	\caption*{#1}
		% \end{figure}

		\fi	
 }


 % hack om na hoofding van definition/proposition/... als dan niet op een nieuwe lijn te starten
% soms is dat goed en mooi, en soms niet; en in HTML is het nu (2/2020) anders dan in pdf
% vandaar suggestie om 
%     \begin{definition}[Nieuw concept] \nl
% te gebruiken als je zeker een newline wil na de hoofdig en titel
% (in het bijzonder itemize zonder \nl is 'lelijk' ...)
\ifdefined\HCode
\newcommand{\nl}{}
\else
\newcommand{\nl}{\ \par} % newline (achter heading van definition etc.)
\fi

% --- samenvatting boc ---
% Dit commando toont een tcolorbox in de PDF en een definition in de HTML!
\newcommand{\samenbox}[2]{
	\ifdefined\HCode %Href werkt niet?
	    \begin{conclusion}[foldable=true,nonumbered,title={#1}]\nl
            #2
        \end{conclusion}
	\else
        \begin{tcolorbox}[title={#1},colback=white,colframe=black,sharp corners]
            #2
        \end{tcolorbox}
		\fi	

}
% ----------------------------------



% % ----------------------------------------------------------------------
% %           User Created Environments --> alles in commentaar gezet, gaan proberen zonder dit te doen
% % ----------------------------------------------------------------------
% % marginnotes voor theorie, samenvatting, bewijs, oefeningen, voorbeeld, oefenmoment.
% % sidenote voor herhaling, oplossing, opmerking en toepassing.
% %
% %kleuren:		theorie - rood
% %			samenvatting - goudgeel
% %			definitie - groen geel
% %			bewijs - blauw
% Oude commands:
% % --- toepassing --> springt niet in enzo...
% \theoremstyle{claim}
% \newtheorem{toepassing}{Toepassing}%[chapter]%<-----------------chapter
% \tcolorboxenvironment{toepassing}{
%   boxrule=0pt,
%   boxsep=0pt,
%   blanker,
%   borderline west={1.5pt}{0pt}{black},
%   sharp corners,
%   before skip=10pt,
%   after skip=10pt,
%   left=5pt,
%   right=5pt,
%   breakable,
% }


% \tcolorboxenvironment{theorem}{
%   boxrule=0pt,
%   boxsep=0pt,
%   colback={White!90!Red},
%   enhanced jigsaw, 
%   borderline west={1pt}{0pt}{Red},
%   sharp corners,
%   before skip=10pt,
%   after skip=10pt,
%   left=5pt,
%   right=5pt,
%   breakable,
% }
% % --- theorie kader
% \theoremstyle{theorem}
% \newtheorem{theorie}{\colorbox{White!90!Red}{Theorie}}%%%[chapter] %<-----------------chapter
% \tcolorboxenvironment{theorie}{
%   boxrule=0pt,
%   boxsep=0pt,
%   colback={White!90!Red},
%   enhanced jigsaw, 
%   borderline west={1.5pt}{0pt}{Red},
%   borderline east={0.5pt}{0pt}{Red},
%   borderline north={0.5pt}{0pt}{Red},
%   borderline south={0.5pt}{0pt}{Red},  
%   sharp corners,
%   before skip=10pt,
%   after skip=10pt,
%   left=5pt,
%   right=5pt,
%   breakable,
% }


% %\declaretheorem[style=theorem]{theorie}

% % \let\proof\relax
% % \let\endproof\relax

% % \declaretheoremstyle[
% %     headfont={\small\scshape},
% %     notefont={\normalfont},
% %     bodyfont={\normalfont},
% %     headpunct={\relax},
% %     headformat={%
% %         \makebox[0pt][r]{\NAME\ \NUMBER\hspace{\marginparsep}}\hskip-\spacelength{\NOTE}},
% % ]{proof}
% % --- Samenvatting
% \theoremstyle{theorem}
% \newtheorem{samenvatting}{\colorbox{White!90!Goldenrod}{Samenvatting}}[chapter]%<-----------------chapter


% \tcolorboxenvironment{samenvatting}{
%   boxrule=0pt,
%   boxsep=0pt,
%   colback={White!90!Goldenrod},
%   enhanced jigsaw, 
%   borderline west={1.5pt}{0pt}{Goldenrod},
%   borderline east={0.5pt}{0pt}{Goldenrod},
%   borderline north={0.5pt}{0pt}{Goldenrod},
%   borderline south={0.5pt}{0pt}{Goldenrod},
%   sharp corners,
%   before skip=10pt,
%   after skip=10pt,
%   left=5pt,
%   right=5pt,
%   breakable,
% }
% % --- definitie
% \theoremstyle{theorem}
% \newtheorem{definitie}{\colorbox{White!90!GreenYellow}{Definitie}}[chapter]%<-----------------chapter


% \tcolorboxenvironment{definitie}{
%   boxrule=0pt,
%   boxsep=0pt,
%   colback={White!90!GreenYellow},
%   enhanced jigsaw, 
%   borderline west={1.5pt}{0pt}{GreenYellow},
%   borderline east={0.5pt}{0pt}{GreenYellow},
%   borderline north={0.5pt}{0pt}{GreenYellow},
%   borderline south={0.5pt}{0pt}{GreenYellow},
%   sharp corners,
%   before skip=10pt,
%   after skip=10pt,
%   left=5pt,
%   right=5pt,
%   breakable,
% }
% % --- definitie
% \theoremstyle{theorem}
% \newtheorem{wet}{\colorbox{White!90!Green}{Wet}}[chapter]%<-----------------chapter


% \tcolorboxenvironment{wet}{
%   boxrule=0pt,
%   boxsep=0pt,
%   colback={White!90!Green},
%   enhanced jigsaw, 
%   borderline west={1.5pt}{0pt}{Green},
%   borderline east={0.5pt}{0pt}{Green},
%   borderline north={0.5pt}{0pt}{Green},
%   borderline south={0.5pt}{0pt}{Green},
%   sharp corners,
%   before skip=10pt,
%   after skip=10pt,
%   left=5pt,
%   right=5pt,
%   breakable,
% }
% % --- bewijs
% \tcolorboxenvironment{bewijs}{
%   boxrule=0pt,
%   boxsep=0pt,
%   blanker,
%   borderline west={1.5pt}{0pt}{NavyBlue!80!white},
%   borderline east={0.5pt}{0pt}{NavyBlue!80!white},
%   borderline north={0.5pt}{0pt}{NavyBlue!80!white},
%   borderline south={0.5pt}{0pt}{NavyBlue!80!white},
%   before skip=10pt,
%   after skip=10pt,
%   left=5pt,
%   right=5pt,
%   breakable,
% }

% \declaretheoremstyle[
%     headfont={\footnotesize\itshape},
%     notefont={\normalfont},
%     bodyfont={\normalfont},
%     headpunct={\relax},
%     headformat={%
%         \makebox[0pt][r]{\NAME\hspace{\marginparsep}}\hskip-\spacelength{\NOTE}},
% ]{claim}

% \declaretheorem[
%     % style=proof,
%     style=claim,
%     qed=\qedsymbol]{bewijs}

% \declaretheorem[style=claim]{Intuition}


% % --- vraag
% \theoremstyle{theorem}
% \newtheorem{vraag}{Vraag}[chapter]%<-----------------chapter
% \tcolorboxenvironment{vraag}{
%   boxrule=0pt,
%   boxsep=0pt,
%   blanker,
% %  borderline west={1pt}{0pt}{Black},
%   borderline south={0.5pt}{-5pt}{Black},
%   %borderline north={1pt}{-3pt}{Black},
%   before skip=10pt,
%   after skip=10pt,
%   sharp corners,
%   left=5pt,
%   right=5pt,
%   breakable,
% }

% % --- Oefening
% \theoremstyle{theorem}
% \newtheorem{oef}{Oefening}[chapter]%<-----------------chapter
% \tcolorboxenvironment{oef}{
%   boxrule=0pt,
%   boxsep=0pt,
%   blanker,
% %  borderline west={1pt}{0pt}{Black},
%   borderline south={0.5pt}{-5pt}{Black},
%   %borderline north={1pt}{-3pt}{Black},
%   before skip=10pt,
%   after skip=10pt,
%   sharp corners,
%   left=5pt,
%   right=5pt,
%   breakable,
% }
% % --- voorbeeld
% \theoremstyle{theorem}
% \newtheorem{voorbeeld}{\colorbox{White!90!Cerulean}{Voorbeeld}}[chapter]%<-----------------chapter
% \tcolorboxenvironment{voorbeeld}{
%   boxrule=0pt,
%   boxsep=0pt,
%   colback={White!90!Cerulean},
%   enhanced jigsaw, 
%   borderline west={1.5pt}{0pt}{Cerulean},
%   borderline north={0.5pt}{0pt}{Cerulean},
%   borderline south={0.5pt}{0pt}{Cerulean},
%   borderline east={0.5pt}{0pt}{Cerulean},
%   sharp corners,
%   before skip=10pt,
%   after skip=15pt,
%   left=5pt,
%   right=5pt,
%   breakable,
% }
% % --- oefenmoment: quizjes tijdens de cursus om te testen of je mee bent. Oplossingen staan in oplossingen
% \theoremstyle{theorem}
% \newtheorem{quiz}{\colorbox{White!90!NavyBlue}{Quiz}}[chapter]%<-----------------chapter
% \tcolorboxenvironment{quiz}{
%   boxrule=0pt,
%   boxsep=0pt,
%   blanker,
%   borderline north={1pt}{-5pt}{NavyBlue!80!white},
%   borderline south={1pt}{-5pt}{NavyBlue!80!white},
%   before skip=10pt,
%   after skip=10pt,
%   sharp corners,
%   left=5pt,
%   right=5pt,
%   %breakable,
% }

% % ----- oplossingen van de quisjes
% \theoremstyle{theorem}
% \newtheorem{oplq}{\colorbox{White!90!NavyBlue}{Opl. Quiz}}[chapter]%<-----------------chapter
% \tcolorboxenvironment{oplq}{
%   boxrule=0pt,
%   boxsep=0pt,
%   blanker,
%   borderline west={1pt}{0pt}{NavyBlue},
%   before skip=10pt,
%   after skip=10pt,
%   sharp corners,
%   left=5pt,
%   right=5pt,
%   breakable,
% }
% % ----- antwoorden op de vragen
% \theoremstyle{theorem}
% \newtheorem{ant}{\colorbox{White!90!black}{Antwoord}}[chapter]%<-----------------chapter
% \tcolorboxenvironment{ant}{
%   boxrule=0pt,
%   boxsep=0pt,
%   blanker,
% %  borderline west={1pt}{0pt}{Black},
%   borderline south={1pt}{-5pt}{Black},
%   %borderline north={1pt}{-3pt}{Black},
%   before skip=10pt,
%   after skip=10pt,
%   sharp corners,
%   left=5pt,
%   right=5pt,
%   breakable,
% }

% % ----- oplossingen van de vraagstukken
% \theoremstyle{theorem}
% \newtheorem{opl}{\colorbox{White!90!black}{Oplossing}}[chapter]%<-----------------chapter
% \tcolorboxenvironment{opl}{
%   boxrule=0pt,
%   boxsep=0pt,
%   blanker,
% %  borderline west={1pt}{0pt}{Black},
%   borderline south={1pt}{-5pt}{Black},
%   %borderline north={1pt}{-3pt}{Black},
%   before skip=10pt,
%   after skip=10pt,
%   sharp corners,
%   left=5pt,
%   right=5pt,
%   breakable,
% }
% % --- herhaling
% \theoremstyle{claim}
% \newtheorem{her}{Herhaling}[chapter]%<-----------------chapter


% % --- oplossingen, gebruik \label en \ref
% %\theoremstyle{claim}
% %\newtheorem{opl}{Oplossing}[chapter]%<-----------------chapter

%  % --- opmerkingen
% \theoremstyle{claim}
% \newtheorem{opm}{Opmerking}[chapter]%<-----------------chapter
% \tcolorboxenvironment{opm}{
%   boxrule=0pt,
%   boxsep=0pt,
%   blanker,
%   borderline west={1.5pt}{0pt}{lightgray},
%   sharp corners,
%   before skip=10pt,
%   after skip=10pt,
%   left=5pt,
%   right=5pt,
%   breakable,
% }

% % --- toepassing
% \theoremstyle{claim}
% \newtheorem{toepassing}{Toepassing}[chapter]%<-----------------chapter
% \tcolorboxenvironment{toepassing}{
%   boxrule=0pt,
%   boxsep=0pt,
%   blanker,
%   borderline west={1.5pt}{0pt}{Black},
%   sharp corners,
%   before skip=10pt,
%   after skip=10pt,
%   left=5pt,
%   right=5pt,
%   breakable,
% }

% \renewcommand*{\raggedleftmarginnote}{\noindent}
% \renewcommand*{\raggedrightmarginnote}{\noindent}
% \newcommand{\mn}[1]{\textsuperscript{\thesidenote}{}\marginnote{\textsuperscript{\thesidenote}{}\itshape\footnotesize #1}\refstepcounter{sidenote}}
% \newcommand{\en}[1]{\marginnote{\footnotesize{#1}}}
% \newcommand{\lec}[2]{{\setlength{\marginparwidth}{.07\paperwidth}\reversemarginpar\marginnote{\centering\footnotesize{\textsf{\bfseries #1}}\\#2}}}
% \newcommand{\sn}[1]{\sidenote{\itshape\footnotesize #1}}

% % -- fullpage
% \newenvironment{fullpage}
%     {\smallskip\noindent\begin{minipage}
%     {\textwidth+\marginparwidth+\marginparsep}\hrule\smallskip\smallskip}
%     {\hrule\end{minipage}\vspace{.1in}}

% % --- varwidth
% \newtcbox{\varbox}{colframe=black,
% colback=white,varwidth upper}

% % voor kolom
% \renewcommand{\arraystretch}{2}
% 	\newcolumntype{P}[1]{>{\centering\arraybackslash}p{#1}}

% % tikz
% \tikzstyle{charge}=[thin,top color=red!50,bottom color=red!70,shading angle=20]
% \tikzstyle{charge+}=[thin,top color=red!50,bottom color=red!90!black,shading angle=20]
% \tikzstyle{charge-}=[thin,top color=blue!50,bottom color=blue!80,shading angle=20]
% \tikzstyle{force}=[->,very thick,orange!80!black]
% \tikzstyle{vector}=[->,very thick,green!45!black]
% \def\a{1.4}
% \def\R{0.15}
% \def\F{0.8}
% %\def\s{1}		%scale van figuren


% ----------------------------------------------------------------------
%           Boek informatie				als je met maketitle werkt
% ----------------------------------------------------------------------
% Information about the book
% to be put on the title page
%\title{Boek template}
%\author{Daan Lipkens}
%\date{Sint-Ursula Instituut \today}




%% Fix 'solution' 
\ifdefined\HCode%
    \renewenvironment{solution}[1][onzichtbaar]
   {%
       \begin{expandable}{xmoplossing}{}\par   %% without \part tikz-error with \begin{oplossing}\tikz{...} !!!!
   }
   {%
   	   \end{expandable}
   } 
\fi
